\documentclass[TFG,MEM,EN,no-cover,no-toc]{P3CTeX}
\usepackage{pxDIA}
\PECTeXconfig{no-plagi=false,
  data={student={Integration test},grade={Software Engineering},
    subj-fullname={Final project},subj-shortname={TFG},
    ca-fullname={Memory},ca-shortname={MEM}}}
\begin{document}
\chapter{Computational diagrams}
\pxTCompiler[caption={Main compiler},label={fig:main-compiler}]{C}{x86}{C}
\pxGraph[caption={Processing graph},label={fig:main-graph}]
  {start / Input / 0,0 / start; end / Output / 2,0 / goal}
  {start->end / 1}
Main compiler: \Cref{fig:main-compiler}. Graph: \ref{fig:main-graph}.
\pxMEMannex[label={ann:diagram}]{Diagram evidence}{
  \pxTCompiler[caption={Compiler in annex},label={fig:annex-compiler}]{Pascal}{x86}{C}
  \pxGraph[caption={Graph in annex},label={fig:annex-graph}]
    {a / A / 0,0; b / B / 2,0}{a->b}
  Annex figure: \Cref{fig:annex-compiler} and \ref{fig:annex-graph}.
}
Annex: \Cref{ann:diagram}.
\end{document}
